\section{Introduction}

Documents are fundamental carriers of human knowledge, encompassing a vast array of formats including academic papers, business reports, legal contracts, medical records, and financial statements. The automatic understanding of these documents---extracting structured information, answering questions, and supporting decision-making---has long been a central challenge in artificial intelligence. This field, broadly termed \textbf{Document Intelligence}, has witnessed remarkable evolution over the past decades, driven by advances in computer vision, natural language processing, and more recently, multimodal deep learning.

The document processing pipeline traditionally follows a two-stage paradigm: \textit{perception} followed by \textit{reasoning}. In the perception stage, documents are first converted into machine-readable representations through Optical Character Recognition (OCR) and layout analysis, identifying text regions, tables, figures, and structural elements. The reasoning stage then operates on these structured representations to perform higher-level tasks such as information extraction, question answering, and knowledge synthesis. However, this modular design introduces several challenges: error propagation from OCR to downstream tasks, loss of visual and spatial information during text linearization, and the inability to handle complex layouts that require joint understanding of text, structure, and visual appearance.

With the emergence of large language models (LLMs) and vision-language models (VLMs), the field is undergoing a paradigm shift toward end-to-end, multimodal document understanding. Modern systems can now process document images directly, bypassing explicit OCR stages while achieving superior performance on complex reasoning tasks. This evolution raises fundamental questions: What are the key technical innovations that enabled this transition? How do different approaches---traditional pipelines, retrieval-augmented generation, and unified multimodal models---compare in terms of capabilities and limitations? What evaluation benchmarks best capture real-world deployment requirements?

\subsection{Scope and Contributions}

This survey provides a comprehensive review of document intelligence research, organized around the core pipeline illustrated in \cref{fig:pipeline_overview}. Our contributions are threefold:

\begin{enumerate}[leftmargin=*]
    \item \textbf{Structured Taxonomy}: We present a systematic organization of the field along two dimensions---perception (layout analysis and OCR) and reasoning (table understanding, text-based RAG, and multimodal VLM approaches). This structure reflects the historical evolution while highlighting emerging unified paradigms.
    
    \item \textbf{Technical Depth}: For each research direction, we trace the evolution from early heuristic methods through deep learning innovations to contemporary LLM/VLM-based approaches, emphasizing key architectural decisions, training strategies, and their impact on performance.
    
    \item \textbf{Evaluation Perspective}: We provide a comprehensive overview of benchmarks spanning single-page to multi-page documents, synthetic to real-world distributions, and accuracy-focused to robustness-oriented evaluation protocols, identifying gaps between research metrics and deployment requirements.
\end{enumerate}

\subsection{Organization}

The remainder of this survey is organized as follows. \cref{sec:layout_ocr} covers Part~I: Document Layout Analysis and OCR, including geometric structure detection, layout-aware representation learning, and the evolution of OCR from traditional methods to deep learning and VLM-based approaches. \cref{sec:table,sec:text_rag,sec:vlm} constitute Part~II: Document Understanding and Question Answering, covering table understanding, text-based retrieval-augmented generation, and vision-language models respectively. \cref{sec:eval} reviews evaluation benchmarks and datasets. Finally, \cref{sec:conclusion} discusses future research directions and concludes the survey.
